%%
%%  Department of Electrical, Electronic and Computer Engineering.
%%  EPR400/2 Final Report - metric definitions of the qualification tests.
%%
%%  Every accuracy, distance and agreement measure quoted in Section 4 is
%%  defined here in full.  Section 4 names each measure, states the reading it
%%  uses and quotes the value; this appendix gives the defining equation, the
%%  routine that implements it and the reason the reading matters.  Nothing
%%  here is a result: no value measured on the system appears in this file.
%%

\subsection{Metric definitions of the qualification tests}
\label{sec:app-met}

The qualification tests of Section~\ref{sec:res-qtp} quote accuracies, correlations and distances whose definitions are not interchangeable: three different averages of character accuracy over a set of lines differ by several points on the same data, and the word accuracy reported throughout is a recall rather than one minus the word error rate.
This appendix defines each measure once, names the routine that computes it, and states where an alternative reading would change the number.
Section~\ref{sec:res-qtp} refers to this appendix in place of repeating the mathematics beside each result.

\subsubsection{Confidence intervals and resampling}

Proportions are reported with a 95\% Wilson score interval~\cite{wilson1927},
\begin{equation}
\label{eq:res-wilson}
\hat p_{\pm}=\frac{\hat p+\dfrac{z^2}{2n}\pm z\sqrt{\dfrac{\hat p(1-\hat p)}{n}+\dfrac{z^2}{4n^2}}}{1+\dfrac{z^2}{n}},\qquad \hat p=\frac{k}{n},
\end{equation}
where $k$ is the number of successes in $n$ trials and $z=1.96$.
Equation~\ref{eq:res-wilson} is used in preference to the normal approximation because it stays valid near 0\% and 100\%, where the system's per-writer rates lie.
Characters of one page are not independent, so accuracies use a page-level percentile bootstrap with 2000 resamples~\cite[p.~3]{efron1979} rather than a closed-form interval, and a requirement is declared met only if the lower confidence bound exceeds the target.

\subsubsection{Character accuracy and its three averages}

Character accuracy of one line uses the Levenshtein distance~\cite[pp.~707--708]{levenshtein1966} $d_L(h,r)$ between hypothesis $h$ and reference $r$, the minimum number of insertions, deletions and substitutions, evaluated by the dynamic-programming recursion of Wagner and Fischer~\cite{wagner1974},
\begin{equation}
\label{eq:res-cacc}
A_{\mathrm{c}}(h,r)=\max\left(0,\;1-\frac{d_L(h,r)}{\lvert r\rvert}\right),
\end{equation}
where $\lvert r\rvert$ is the reference length.
Equation~\ref{eq:res-cacc} is implemented as \texttt{Classify.CharAcc}, which floors the value at zero so that a hypothesis longer than the reference cannot score negatively.

A set of lines can be summarised in three different ways, and they do not agree, so each is named explicitly:
\begin{equation}
\label{eq:res-cacc-avg}
\bar A_{\mathrm{c}}^{\,\mathrm{line}}=\frac{1}{L}\sum_{\ell=1}^{L}A_{\mathrm{c}}(h_\ell,r_\ell),\qquad
\bar A_{\mathrm{c}}^{\,\mathrm{micro}}=1-\frac{\sum_\ell d_L(h_\ell,r_\ell)}{\sum_\ell\lvert r_\ell\rvert},\qquad
\bar A_{\mathrm{c}}^{\,\mathrm{writer}}=\frac{1}{W}\sum_{w=1}^{W}\bar A_{\mathrm{c}}^{\,\mathrm{line}}(w),
\end{equation}
the unweighted mean over lines, the character-weighted micro-average over all lines pooled, and the unweighted mean over the $W$ writers of their own line means.
The micro-average of Equation~\ref{eq:res-cacc-avg} is what ``summed over all sentences'' would mean, but the character-accuracy figures quoted for qualification test~1 are the \emph{mean over writers}, $\bar A_{\mathrm{c}}^{\,\mathrm{writer}}$, recorded as \texttt{char\_acc\_mean\_over\_writers}.
The distinction is not cosmetic: for the held-out judge on the 123 real lines the two readings differ by 4.7 points, 86.6\% over writers against 81.9\% over lines.
Figures quoted for qualification test~2 are instead the unweighted mean over line crops, $\bar A_{\mathrm{c}}^{\,\mathrm{line}}$, because \texttt{verify\_text\_accuracy.py} accumulates one \texttt{CharAcc} per crop and divides by the crop count, so a two-word crop weighs as much as a twenty-word one.

\subsubsection{Word accuracy is a recall, not one minus the word error rate}

Word accuracy, which Section~\ref{sec:res-qtp} quotes six times, is \texttt{Classify.WordAcc} and is \emph{not} $1-\mathrm{WER}$:
\begin{equation}
\label{eq:res-wacc}
A_{\mathrm{w}}^{\,\mathrm{word}}=\frac{\mathrm{LCS}(\mathrm{words}(h),\mathrm{words}(r))}{\lvert\mathrm{words}(r)\rvert},
\end{equation}
the length of the longest common subsequence of the two word sequences, compared case-insensitively, divided by the number of reference words.
Equation~\ref{eq:res-wacc} is therefore an order-preserving \emph{recall} over words: it counts how many of the intended words were recovered in the right order and imposes no penalty for inserted or spurious words, whereas $1-\mathrm{WER}$ charges insertions, deletions and substitutions alike and can go negative.
A reader who assumes $1-\mathrm{WER}$ will read every word figure in Section~\ref{sec:res-qtp} as stricter than it is.

\subsubsection{Writer accuracy is reported per line, not per page}

Writer identification admits a per-line and a per-page reading,
\begin{equation}
\label{eq:res-wacc-writer}
A_{\mathrm{w}}^{\,\mathrm{line}}=\frac{1}{L}\sum_{\ell=1}^{L}\mathbf{1}\!\left[\hat w_\ell=w_\ell\right],\qquad
A_{\mathrm{w}}^{\,\mathrm{page}}=\frac{1}{P}\sum_{p=1}^{P}\mathbf{1}\!\left[\arg\max_w \tfrac{1}{L_p}\textstyle\sum_{\ell\in p}P(w\mid\ell)=w_p\right],
\end{equation}
the fraction of lines assigned to the correct writer, and the fraction of pages whose mean per-line writer probability peaks at the correct writer.
The proposal's wording implies the page form of Equation~\ref{eq:res-wacc-writer}, but $A_{\mathrm{w}}^{\,\mathrm{page}}$ was never evaluated: every value reported in Section~\ref{sec:res-qtp}, including the 99.78\% and 93.67\% over 1358 lines, is the per-line form $A_{\mathrm{w}}^{\,\mathrm{line}}$.
Because lines of one page share a writer, the page form would be the stronger claim and the smaller sample.
\mustdo{evaluate $A_{\mathrm{w}}^{\,\mathrm{page}}$ of Equation~\ref{eq:res-wacc-writer} by pooling the stored per-line probabilities of \texttt{stage4\_predictions\_ink.json} and \texttt{stage4\_predictions\_shape.json} by page.}

\subsubsection{Decoding}

All text figures in the report were produced by greedy connectionist temporal classification decoding --- argmax per time step, collapse repeats, drop the blank, as implemented in \texttt{Classify.DecodeCtc} --- with \textbf{no beam search, no language model and no lexicon} anywhere in the deployed pipeline.
That matters next to the 94.21\% word accuracy of qualification test~2: the figure is what an unaided per-frame argmax recovers, so it is not inflated by a decoder that could snap a near-miss onto a dictionary word; equally, nothing corrects a single-character slip that destroys a word.
All random choices use recorded seeds.

\subsubsection{Inter-rater agreement of the panel that was not convened}

Inter-rater agreement of the blind human panel of qualification test~1 was to be quantified with Fleiss' kappa~\cite{fleiss1971},
\begin{equation}
\label{eq:res-kappa}
\kappa=\frac{\bar P-\bar P_e}{1-\bar P_e},
\end{equation}
where $\bar P$ is the mean observed pairwise agreement and $\bar P_e$ the agreement expected by chance.
As the panel was never convened, Equation~\ref{eq:res-kappa} was never evaluated and is given here only to define the sub-test that was not run; the worksheet prepared for it is Table~\ref{tab:res-panel}.

\subsubsection{Style correlation on geometric features}

Each feature is $z$-scored, $z_{wf}=(v_{wf}-\mu_f)/\sigma_f$, using the mean and standard deviation over the ten real writers, and the stylistic correlation reported for sub-test~(b) of qualification test~1 is the Pearson product-moment correlation coefficient~\cite{pearson1896},
\begin{equation}
\label{eq:res-pearson}
\rho=\frac{\sum_{i}(a_i-\bar a)(b_i-\bar b)}{\sqrt{\sum_i(a_i-\bar a)^2}\sqrt{\sum_i(b_i-\bar b)^2}},
\end{equation}
where $a_i$ and $b_i$ are the real and synthesised $z$-scored entries pooled over writers and features (70 pairs for the seven-feature set the protocol prescribes, listed in Table~\ref{tab:res-features}).
Besides the pooled value of Equation~\ref{eq:res-pearson}, a per-writer value (7 pairs) and a nearest-centroid retrieval in $z$-space with its $10\times10$ confusion matrix are also computed.
The correlation actually measured used the nine geometric features of Table~\ref{tab:res-r1} rather than the seven of Table~\ref{tab:res-features}.
\mustdo{confirm the Pearson reference: the product-moment correlation coefficient is introduced in the 1896 paper ``Regression, heredity, and panmixia'' (Philosophical Transactions of the Royal Society A, vol.~187), whereas the sourcing audit suggested vol.~186 (1895), which is the skew-variation paper. Verify the volume, year and page range against the original and correct the bibliography entry.}

\subsubsection{Overlay distances in the pixel domain}

Both images entering the overlay test are normalised in the same way before any distance is taken, so that the comparison measures shape rather than size or placement; Figure~\ref{fig:res-overlay} sets out that procedure.

\begin{figure}[h]
\centering
\begin{tikzpicture}[
  resbox/.style={draw,rounded corners=2pt,fill=gray!12,align=center,font=\footnotesize,minimum height=7mm,inner sep=3pt},
  resdark/.style={draw,rounded corners=2pt,fill=gray!30,align=center,font=\footnotesize,minimum height=7mm,inner sep=3pt},
  resarr/.style={-{Stealth[length=2.2mm]},thick}]
\node[resbox,text width=2.9cm] (real) at (0,0.9) {Genuine sentence image (held-out form, writer $w$)};
\node[resbox,text width=2.9cm] (syn) at (0,-0.9) {Synthesised image of the same text in style $w$ (or $w'$)};
\node[resbox,text width=2.7cm] (n1) at (4.1,0.9) {Crop to ink, scale to $x_{\mathrm{h}}$, align centroid};
\node[resbox,text width=2.7cm] (n2) at (4.1,-0.9) {Crop to ink, scale to $x_{\mathrm{h}}$, align centroid};
\node[resdark,text width=2.4cm] (d) at (8.0,0) {Skeleton point sets $A$, $B$: $d_{\mathrm{ch}}$, $d_{\mathrm{MH}}$};
\node[resbox,text width=2.3cm] (out) at (11.4,0) {Matched vs mismatched: AUC, rank-1};
\draw[resarr] (real)--(n1); \draw[resarr] (syn)--(n2);
\draw[resarr] (n1.east)--(d.north west); \draw[resarr] (n2.east)--(d.south west);
\draw[resarr] (d)--(out);
\end{tikzpicture}
\caption{Overlay test of qualification test 1: both images are normalised in the same way and matched writer pairs are compared with mismatched pairs.}
\label{fig:res-overlay}
\end{figure}

The skeleton point sets $A$ and $B$ that this procedure produces are compared with
\begin{equation}
\label{eq:res-chamfer}
d_{\mathrm{ch}}(A,B)=\frac{1}{2}\left[\frac{1}{\lvert A\rvert}\sum_{a\in A}\min_{b\in B}\lVert a-b\rVert+\frac{1}{\lvert B\rvert}\sum_{b\in B}\min_{a\in A}\lVert a-b\rVert\right],
\end{equation}
\begin{equation}
\label{eq:res-mhd}
d_{\mathrm{MH}}(A,B)=\max\left[\frac{1}{\lvert A\rvert}\sum_{a\in A}\min_{b\in B}\lVert a-b\rVert,\;\frac{1}{\lvert B\rvert}\sum_{b\in B}\min_{a\in A}\lVert a-b\rVert\right],
\end{equation}
\begin{equation}
\label{eq:res-auc}
\mathrm{AUC}=\Pr\left[d_{\mathrm{matched}}<d_{\mathrm{mismatched}}\right]+\tfrac{1}{2}\Pr\left[d_{\mathrm{matched}}=d_{\mathrm{mismatched}}\right],
\end{equation}
where $\lVert\cdot\rVert$ is the Euclidean norm and distances are in units of the x-height $x_{\mathrm{h}}$.
Equation~\ref{eq:res-chamfer} is the mean of the two directed mean nearest-neighbour distances, which symmetrises the one-directional chamfer matching distance of Barrow \emph{et al.}~\cite[pp.~659--660]{barrow1977}; Equation~\ref{eq:res-mhd} is their maximum, which is less sensitive to outliers than the classical Hausdorff distance~\cite{dubuisson1994}; and Equation~\ref{eq:res-auc} is the area under the receiver operating characteristic curve~\cite[p.~868]{fawcett2006}, 0.5 for a distance carrying no writer information and 1.0 for perfect separation.


The protocol sized this sub-test at 8 writers and 5 sentences, which gives 40 matched and $40\times9=360$ mismatched pairs; the run actually made used ten writers and three sentences each, which is the 30 matched and 300 mismatched pairs reported in Section~\ref{sec:res-qtp1}.
The features of Table~\ref{tab:res-features} are computed by one deterministic extractor with identical thresholds on every image, after cropping to the ink and binarising with the same rule.

\subsubsection{Positional error and its uncertainty budget}

The absolute error of a drawn cross and the uncertainty of its measurement, prescribed for qualification test~5 and never evaluated, are
\begin{equation}
\label{eq:res-r5}
e_i=\lVert\mathbf{p}_i-\mathbf{p}_i^{\mathrm{cmd}}\rVert,
\end{equation}
\begin{equation}
\label{eq:res-unc}
u=\sqrt{u_{\mathrm{scan}}^2+u_{\mathrm{ref}}^2+u_{\mathrm{id}}^2},
\end{equation}
where $\mathbf{p}_i$ is the measured and $\mathbf{p}_i^{\mathrm{cmd}}$ the commanded centre of cross $i$, referenced to the machine origin 5~mm inside the minimum end-stops.
In Equation~\ref{eq:res-unc}, $u_{\mathrm{scan}}$ is the scanner scale error obtained by scanning the verified ruler, $u_{\mathrm{ref}}$ the registration of the sheet against the machine origin and $u_{\mathrm{id}}$ the error of locating a cross centre, taken as half a pixel; the three standard uncertainties are combined in quadrature and the expanded uncertainty is $2u$, following the Guide to the Expression of Uncertainty in Measurement~\cite[Clause~5]{jcgm2008}.
Mean, maximum and 95th percentile of the error of Equation~\ref{eq:res-r5} were to be reported together with the uncertainty of Equation~\ref{eq:res-unc}, on the sheet reproduced blank as Table~\ref{tab:res-sheet}.
Because no mark was ever drawn, neither equation was evaluated.
